\chapter{Conclusions}%
\label{chp:conclusions}

The purpose of this final year project was to design, implement, and evaluate a modular computer vision pipeline capable of producing offside verdicts from broadcast football footage, without requiring the dedicated multi-camera infrastructure used by professional \gls{var} systems. The seven objectives stated in Section~\ref{sec:intro_aims} have each been addressed and fulfilled, demonstrating that the core aim is indeed achievable.

\section{Review of Objectives}%
\label{sec:conc_objectives}

The first three objectives focused on data acquisition and preparation. Matches and associated ground-truth annotations were gathered from SoccerNet across six competitions and multiple seasons. A total of 457 offside clips were extracted and filtered, each trimmed to a 35-second window around the annotated event. A browser-based keyframe annotation interface allowed the operator to identify the exact frame at which the ball was played; out of the 457 clips, 164 yielded usable keyframes and 293 were treated as unsuitable.

The fourth objective was player detection and team assignment. The keyframes were analysed using the YOLO11x-seg algorithm. On average 18.8 players were detected per frame, with no frame producing zero detections. The team was assigned using a 6D CIE Lab feature vector which was clustered using K-Means, successfully assigning the teams in most of the 164 cases. Team misidentification, attributable to spectrally similar kits, accounted for three of the 22 disagreements on the clean evaluation set; the remaining 19 disagreements were marginal or moderate cases within the geometric uncertainty of the calibration.

The fifth objective was to estimate a homography. A two-pass RANSAC with least-squares refit achieved a mean symmetric pixel reprojection error of 1.00\,px across the 150 retained calibrated frames (85.3\% below 2\,px), corresponding to a world-space uncertainty of approximately 15--25\,cm.

The sixth objective was to evaluate the offside condition geometrically. On the 150 clean evaluation events, overall agreement with SoccerNet annotations was 85.3\%; for the 92 confident events (margin $>$0.5\,m), agreement was \textbf{91.3\%}, and for clear cases (margin $>$1.0\,m) agreement was 97.4\%.

The seventh objective was to evaluate overall accuracy and identify failure modes. The analysis classified the failures into three categories: frames excluded by the homography quality filter (5 frames with reprojection error exceeding 3.2\,px), team misidentification (three cases), and marginal or moderate measurement uncertainty (19 cases with margin $\leq$1.0\,m).

\section{Principal Conclusions}%
\label{sec:conc_principal}

In conclusion, a system that relies entirely on broadcast footage for offside detection is practical at acceptable levels of accuracy. The almost complete agreement on the clear cases demonstrates that the geometric pipeline is fundamentally correct. In cases where the margin is clearly significant enough, the system is almost guaranteed to make the right call. The drop in agreement at sub-0.5\,m cases is not necessarily a flaw of the algorithm, but rather a mix of both limitations due to a single camera broadcast, and also cases in which the player was actually just slightly onside but the referee flagged it as offside due to the tight margin between the attacker and defender, to which linesmen are prone. 

The system compares favourably with the only comparable end-to-end published result. Uchida \textit{et~al.}\ reported an F1 score of 63.1\% using two fixed calibrated cameras on 20 clips~\cite{uchida2021offside}; the present system achieves a 91.3\% agreement rate on confident cases from a single broadcast feed across a substantially larger evaluation set.

The major challenges to the application of the technique remain: keyframe selection, pitch landmark picking, and the offside checker interaction. These steps are critical to the accuracy of the verdict. An incorrect keyframe directly invalidates the result and their automation should be the main focus of future development. The geometric and detection components of the pipeline are sufficiently accurate that, were these manual steps automated, a fully deployable system would be within reach.

\section{Closing Remarks}%
\label{sec:conc_closing}

The work presented in this final year project demonstrates that the gap between the precision available to elite football through \gls{var} and the information available from a standard broadcast feed is narrower than it might appear. A system built from publicly available tools, operating on footage that already exists for the vast majority of competitive matches, can produce geometrically grounded offside verdicts with high reliability for cases where the margin is large. Closing the remaining gap, automating the manual stages, integrating body keypoint precision, and validating against geometric ground truth are all engineering problems that this work has laid the foundation for.
